
@@M0@@

Si on note @@M1@@ les opérations @@M2@@ correspondant aux opérations @@M3@@, on a donc :
@@M4@@

On notera @@M5@@ le coefficient de @@M6@@ dans @@M7@@~; le lecteur pourra vérifier que c'est un polynôme de @@M8@@, dont les composantes isobares sont de poids variant entre @@M9@@ et @@M10@@.

Soient @@M11@@ un @@M12@@-anneau, @@M13@@. Du diagramme (3.6.1), on déduit les relations :
@@M14@@
@@M15@@

On remarquera enfin que la structure de pré-@@M16@@-anneau qui vient d'être définie sur @@M17@@ est une structure de @@M18@@-anneau, puisque déduite de celle qui a été définie en 2.3 par l'isomorphisme @@M19@@.

3.7. Pour éviter les confusions, le foncteur @@M20@@, muni de la structure de foncteur en @@M21@@-anneaux définie au §2 (resp. au §3) sera noté @@M22@@ (resp. @@M23@@). Il résulte des considérations précédentes que pour se donner une
